\section*{Abstract}
		In the present work, the fundamental architecture of spacetime is reinterpreted within the framework of \textbf{Fundamental Fractal Geometric Field Theory (FFGFT)} – internally referred to as the T0 model (B18). The central paradigm consists in the transition from a point-like to a purely geometric description of the vacuum as a four-dimensional \textbf{Gyral Torus}.
		
		\textbf{Geometric Structure:} The theory is based on the fractal-geometric foundation with the parameter $\xi \approx (4/3)\times 10^{-4}$ and the densest local sphere packing by regular \textbf{Tetrahedra}. This tetrahedral basis forms the stable foundation for the low generations (electron, muon, proton/neutron) as well as the local 3D crystal structure of the torus. Building upon this, the ideal sub-Planck factor
		\begin{equation*}
			f = 7500,
		\end{equation*}
		emerges through fractal branching and pentagonal symmetry breaking, representing an exactly 7500-fold reduction compared to the conventional Planck scale ($t_0$) and following directly from the geometric winding density $30000/4$.
		
		\textbf{g-2 Anomaly:} A core element of the work is the transparent geometric derivation of the anomalous magnetic moments of leptons. While the Standard Model relies on numerous perturbative terms, in FFGFT the electron anomaly follows directly from the base winding (tetrahedral projection). The muon and tau anomalies arise from fractal branchings with Hausdorff dimensions $p \approx 5/3$ and $4/3$, respectively. With the ideal value $f = 7500$ and the geometric projection factor $k_\text{geom} = (2/\sqrt{\varphi})\sqrt{2} = 2.224$, the purely geometric predictions achieve an accuracy of about 2\,\%. Through the \textbf{fractal correction} $K_\text{frak}^{3/2}$, this systematic deviation is explained purely geometrically: the effective projection factor $k_\text{eff} = k_\text{geom}/K_\text{frak}^{3/2} = 2.2692$ agrees with the experimentally reconstructed value $k_\text{rec} = 2.2696$ to 0.01\,\%. The corrected absolute values achieve 0.01\% (electron) and 0.15\% (muon) accuracy. The most precise, $k_\text{geom}$-independent prediction for the tau anomaly is
		\begin{equation*}
			a_\tau \approx 1.277 \times 10^{-3},
		\end{equation*}
		which follows exclusively from the exact ratio $f^{1/3} - 1$.
		
		\textbf{Geometric Proportionality:} All physical base quantities (constants, masses, couplings) stand in fixed geometric ratios, drastically reducing the number of free parameters compared to the Standard Model. The T0 theory thus offers a complete geometric description and provides concrete, experimentally testable predictions – particularly for the tau anomaly as a decisive test at Belle II.

	
	\begin{tcolorbox}[colback=yellow!10!white, colframe=orange!75!black, title=Note on Older Documents]
		Previous versions of the g-2 analysis (018\_T0\_Anomale-g2-9/10) used semi-empirical factors or the ideal $k_\text{geom}$ without fractal correction. The present Rev.~11 shows that the ~2\,\% deviation of absolute values is \textbf{fully explained by the fractal correction $K_\text{frak}^{3/2}$} -- the effective projection factor $k_\text{eff} = k_\text{geom}/K_\text{frak}^{3/2}$ agrees with the experimentally reconstructed value to 0.01\,\%.
	\end{tcolorbox}
	
	\textbf{Keywords:} Anomalous magnetic moment, g-2, T0 theory, Time-Mass Duality, Torsion lattice, Ratio predictions, Koide formula
	
	
	\section{Introduction: Geometric vs. Semi-Empirical Approaches}
	
	\subsection{The Philosophy of T0 Theory}
	
	The T0 theory is based on the principle that \textbf{all} physical constants should follow from the geometric structure of a 4-dimensional torsion lattice. For anomalous magnetic moments this means:
	
	\begin{itemize}
		\item \textbf{NO} hidden fit parameters
		\item \textbf{ONLY} geometric factors: $\varphi$, $\xi$, $f$
		\item Honesty about precision limits
		\item Consistency with other predictions
	\end{itemize}
	
	\subsection{Consistency with Mass Predictions}
	
	The T0 theory predicts lepton masses with ~0.4--1.2\% deviation:
	
	\begin{table}[h]
		\centering
		\resizebox{\textwidth}{!}{\begin{tabular}{lccc}
			\toprule
			\textbf{Lepton} & \textbf{T0 [MeV]} & \textbf{Exp [MeV]} & \textbf{Deviation} \\
			\midrule
			Electron & 0.505 & 0.511 & 1.18\% \\
			Muon & 105.0 & 105.7 & 0.66\% \\
			Tau & 1783 & 1777 & 0.37\% \\
			\bottomrule
		\end{tabular}}
		\caption{Lepton masses in T0}
	\end{table}
	
	\textbf{Expectation:} g-2 should have similar precision (~1\%).
	
	With the fractal correction $K_\text{frak}^{3/2}$, the g-2 predictions achieve 0.01--0.15\% accuracy, comparable to the mass precision.
	
	\section{Physical Fundamentals}
	
	\subsection{What is the Anomalous Magnetic Moment?}
	
	The magnetic moment of a charged spin-1/2 particle is:
	\begin{equation}
		\mu = g \cdot \frac{e}{2m} \cdot \frac{\hbar}{2}
	\end{equation}
	
	where $g$ is the gyromagnetic factor (g-factor).
	
	\textbf{Dirac Prediction:} For a point-like particle: $g = 2$
	
	\textbf{Quantum Effects:} Vacuum polarization, vertex corrections $\Rightarrow g \neq 2$
	
	\textbf{Anomaly:} $a = (g-2)/2$
	
	\textbf{QED Expectation:} $a \approx \alpha/(2\pi) + \mathcal{O}(\alpha^2) \approx 0.00116$
	
	\subsection{T0 Interpretation: Windings in the Torsion Lattice}
	
	In T0 theory, leptons are \textbf{winding structures} in the 4D torsion lattice:
	
	\begin{itemize}
		\item \textbf{Electron:} Simple winding (1st generation)
		\item \textbf{Muon:} Winding with fractal branching (2nd generation)
		\item \textbf{Tau:} More complex fractal structure (3rd generation)
	\end{itemize}
	
	The anomalous moment arises from:
	\begin{enumerate}
		\item The \textbf{rotation} of the winding (spin)
		\item The \textbf{charge distribution} on the winding
		\item The \textbf{projection} 4D $\to$ 3D
	\end{enumerate}
	
	$\Rightarrow$ \textbf{No} point-like charge $\Rightarrow$ $a \neq 0$
	
	\section{Geometric Formulas}
	
	\subsection{Fundamental Parameters}
	
	The T0 theory uses exclusively three geometric fundamental constants:
	
	\begin{align}
		\varphi &= \frac{1 + \sqrt{5}}{2} = 1.618\ldots \quad \text{(Golden Ratio)} \\
		\xi &= \frac{4}{3} \times 10^{-4} = 1.333 \times 10^{-4} \quad \text{(Torsion constant)} \\
		f &= 7500 \quad \text{(Sub-Planck factor)}
	\end{align}
	
	\subsection{The Real Sub-Planck Factor: \(f = 7500\)}
	Now we put everything together: The ideal crystal remains intact, the symmetry breaking only affects the projection factors:
	\begin{equation}
		\boxed{f = 7500}
	\end{equation}
	This is the \textbf{most fundamental number of the T0 theory}. It appears in almost all formulas and describes:
	\begin{itemize}
		\item The number of Sub-Planck cells per Planck length
		\item The density of the torsion lattice
		\item The fundamental frequency of all geometric resonances
	\end{itemize}
	\subsection{The Symmetry Breaking: The Role of the Golden Ratio}
	A perfect, ideal crystal would be completely symmetric. Yet our world shows symmetry breaking on all levels:
	\begin{itemize}
		\item Matter dominates over antimatter
		\item The weak interaction violates parity symmetry
		\item The neutron is heavier than the proton
		\item The three lepton generations have different masses
	\end{itemize}
	In the T0 theory, all these symmetry breakings have a single, geometric origin: the pentagonal symmetry of the crystal, embodied by the \textbf{golden ratio} \(\varphi\).
	The golden ratio \(\varphi = (1+\sqrt{5})/2 = 1.618033989\ldots\) is the irrational number describing pentagonal symmetry. In a perfect pentagon, \(\varphi\) appears everywhere: The ratio of diagonal to side is exactly \(\varphi\).
	Why pentagonal symmetry specifically? For deep mathematical reasons, pentagonal symmetry is the first one that \textbf{cannot tile the plane periodically}. This leads to \emph{quasicrystals} – structures that are ordered but not periodic. Exactly such a quasicrystalline structure is postulated by T0 theory for the Sub-Planck scale.
	The symmetry breaking is quantified in the theory not by a direct subtraction of \(5\varphi\) from the ideal anchor number 7500. Instead, it manifests in the \textbf{fractal correction $K_\text{frak}^{3/2}$}, which corrects the ideal geometric projection factor \(k_\text{geom}\) to the effective value \(k_\text{eff}\):
	\begin{equation}
		k_\text{geom} = \frac{2}{\sqrt{\varphi}} \times \sqrt{2} \approx 2.22357,
	\end{equation}
	which projects the 4D torsion onto the 3D world. The \textbf{fractal correction} yields the effective projection factor:
	\begin{equation}
		k_\text{eff} = \frac{k_\text{geom}}{K_\text{frak}^{3/2}} = \frac{2.224}{0.98007} = 2.2692
	\end{equation}
	The exponent $3/2$ corresponds to the half spatial dimension ($D/2 = 3/2$) and describes how the fractal structure $K_\text{frak} = 1 - 100\xi = 0.98667$ acts across three spatial dimensions on the projection.
	
	\textbf{Crucially:} The experimentally reconstructed value \(k_\text{geom}^\text{rec} = 2.26955\) agrees with $k_\text{eff} = 2.2692$ to \textbf{0.014\%} -- the ~2\% deviation of the ideal formula is not a limitation but is fully explained by the fractal correction.
	\begin{quotation}
		From the ideal 7500 remained the ideal 7500. This number became the new fundamental constant of the universe. It determined how densely the lattice was packed, how quickly torsion could propagate, which resonances were possible.
		Everything we observe today – every particle mass, every force strength, every cosmological constant – is a consequence of this single geometric story: From perfect crystal to pentagonally broken reality, with the breaking quantitatively captured by $K_\text{frak}^{3/2}$.
	\end{quotation}
	
	\subsection{Electron: Base Winding}
	
	\textbf{Formula:}
	\begin{equation}
		a_e = \frac{S_3/f}{k_{\text{geom}}}
		\label{eq:ae}
	\end{equation}
	
	where:
	\begin{itemize}
		\item $S_3 = 2\pi^2 = 19.739$: 3D surface of the 4D winding
		\item $f = 7500$: Sub-Planck scaling
		\item $k_{\text{geom}}$: Geometric projection factor
	\end{itemize}
	
	\textbf{Geometric Projection Factor:}
	\begin{equation}
		k_{\text{geom}} = \frac{2}{\sqrt{\varphi}} \times \sqrt{2}
		\label{eq:kgeom}
	\end{equation}
	
	\textbf{Explanation of Factors:}
	\begin{itemize}
		\item $2/\sqrt{\varphi} = 1.572$: Pentagonal projection (from $\xi$-structure)
		\item $\sqrt{2} = 1.414$: Diagonal projection 4D $\to$ 3D
		\item $k_{\text{geom}} = 2.224$: Completely geometric!
	\end{itemize}
	
	\textbf{Numerical Calculation:}
	\begin{align}
		k_{\text{geom}} &= \frac{2}{\sqrt{1.618}} \times \sqrt{2} = 2.224 \\
		a_e &= \frac{19.739 / 7500}{2.224} \\
		a_e &= 1.184 \times 10^{-3}
	\end{align}
	
	\textbf{Comparison (ideal zeroth order):}
	\begin{itemize}
		\item T0 (ideal): $a_e = 1.184 \times 10^{-3}$
		\item Experiment: $a_e = 1.160 \times 10^{-3}$
		\item Deviation: \textbf{2.03\%}
	\end{itemize}
	
	\textbf{With fractal correction $K_\text{frak}^{3/2}$:}
	\begin{align}
		k_\text{eff} &= \frac{k_\text{geom}}{K_\text{frak}^{3/2}} = \frac{2.224}{0.98007} = 2.2692 \\
		a_e^\text{(corr)} &= \frac{S_3/f}{k_\text{eff}} = \frac{19.739/7500}{2.2692} = 1.1598 \times 10^{-3}
	\end{align}
	\begin{itemize}
		\item T0 (corrected): $a_e = 1.1598 \times 10^{-3}$
		\item Experiment: $a_e = 1.1597 \times 10^{-3}$
		\item Deviation: \textbf{0.014\%}
	\end{itemize}
	
	\subsection{Muon: Fractal Additional Winding}
	
	\textbf{Formula:}
	\begin{equation}
		a_\mu = a_e + \Delta a_{\text{fractal}}
		\label{eq:amu}
	\end{equation}
	
	with
	\begin{equation}
		\Delta a_{\text{fractal}} = \frac{4\pi}{f^{p_\mu}}
		\label{eq:delta_mu}
	\end{equation}
	
	where:
	\begin{itemize}
		\item $p_\mu = 5/3$: Fractal Hausdorff dimension
		\item $4\pi$: Complete torsion revolution
	\end{itemize}
	
	\textbf{Meaning of $p_\mu = 5/3$:}
	
	This is the well-known Hausdorff dimension of:
	\begin{itemize}
		\item Brownian motion in 2D
		\item Self-avoiding random walk
		\item Koch curve (fractal)
	\end{itemize}
	
	$\Rightarrow$ Physically plausible for ``partially branched winding''!
	
	\textbf{Numerical Calculation:}
	\begin{align}
		\Delta a_{\text{fractal}} &= \frac{4\pi}{7500^{5/3}} = 4.373 \times 10^{-6} \\
		a_\mu &= 1.184 \times 10^{-3} + 4.373 \times 10^{-6} \\
		a_\mu &= 1.188 \times 10^{-3}
	\end{align}
	
	\textbf{Comparison (ideal zeroth order):}
	\begin{itemize}
		\item T0 (ideal): $a_\mu = 1.188 \times 10^{-3}$
		\item Experiment: $a_\mu = 1.166 \times 10^{-3}$
		\item Deviation: \textbf{1.89\%}
	\end{itemize}
	
	\textbf{With fractal correction:}
	\begin{align}
		a_\mu^\text{(corr)} &= a_e^\text{(corr)} + \Delta a_\text{fractal} = 1.1598 \times 10^{-3} + 4.373 \times 10^{-6} \\
		a_\mu^\text{(corr)} &= 1.1642 \times 10^{-3}
	\end{align}
	\begin{itemize}
		\item T0 (corrected): $a_\mu = 1.1642 \times 10^{-3}$
		\item Experiment: $a_\mu = 1.1659 \times 10^{-3}$
		\item Deviation: \textbf{0.15\%}
	\end{itemize}
	
	\subsection{Tau: More Complex Fractal Structure}
	
	\textbf{Formula:}
	\begin{equation}
		a_\tau = a_e + \frac{4\pi}{f^{p_\tau}}
		\label{eq:atau}
	\end{equation}
	
	where:
	\begin{itemize}
		\item $p_\tau = 4/3$: Stronger fractal branching
	\end{itemize}
	
	\textbf{Meaning of $p_\tau = 4/3$:}
	
	This is the box-counting dimension of many fractals (e.g., Koch curve, Mandelbrot set).
	
	\textbf{Numerical Calculation:}
	\begin{align}
		\Delta a_{\text{fractal}} &= \frac{4\pi}{7500^{4/3}} = 8.560 \times 10^{-5} \\
		a_\tau &= 1.184 \times 10^{-3} + 8.560 \times 10^{-5} \\
		a_\tau &= 1.269 \times 10^{-3}
	\end{align}
	
	\textbf{Status:} This is a \textbf{prediction} – tau-g-2 has not been measured yet!
	
	\textbf{Corrected prediction:}
	\begin{align}
		a_\tau^\text{(corr)} &= a_e^\text{(corr)} + \Delta a_\text{fractal} = 1.1598 \times 10^{-3} + 8.560 \times 10^{-5} \\
		a_\tau^\text{(corr)} &= 1.2454 \times 10^{-3}
	\end{align}
	
	\section{Two Classes of Predictions: Absolute Values vs. Ratios}
	
	\subsection{Why ~2\% Deviation for Absolute Values?}
	
	\subsection{Where do the ~2\% in the ideal formula come from?}
	
	The ~2\% deviation of the ideal formula (with $k_\text{geom} = 2.224$) is \textbf{no longer an open question} but is fully explained by the fractal correction:
	
	\begin{equation}
		k_\text{eff} = \frac{k_\text{geom}}{K_\text{frak}^{3/2}} = \frac{2.224}{(1-100\xi)^{3/2}} = \frac{2.224}{0.98007} = 2.2692
	\end{equation}
	
	\textbf{Physical meaning:}
	\begin{itemize}
		\item $K_\text{frak} = 1 - 100\xi = 0.98667$ is the universal fractal correction of T0 theory
		\item The exponent $3/2 = D/2$ (half spatial dimension) describes the effect of fractal structure on the 4D→3D projection
		\item This correction is \textbf{not a fit}: $K_\text{frak}$ follows directly from $\xi$ and already appears in the mass formulas
		\item The experimentally reconstructed value $k_\text{rec} = 2.2696$ confirms the correction to 0.014\%
	\end{itemize}
	
	\textbf{Consistency:} The same fractal correction $K_\text{frak}$ appears in:
	\begin{enumerate}
		\item Particle masses: $m_i = r_i \cdot \xi^{p_i} \cdot v \cdot K_\text{frak}$
		\item Gravitational constant: $G = f(\xi, K_\text{frak})$
		\item g-2 anomalies: $k_\text{eff} = k_\text{geom} / K_\text{frak}^{3/2}$
	\end{enumerate}
	
	Thus T0 theory is fully consistent: \textbf{one} correction parameter $K_\text{frak}$ with \textbf{different} exponents depending on the physical context.
	
	\subsection{Ratios are Mathematically Exact}
	
	In contrast to absolute values, \textbf{ratios of differences} are structurally exact:
	
	\begin{equation}
		\frac{\Delta a(\tau - \mu)}{\Delta a(\mu - e)} = \frac{4\pi/f^{4/3} - 4\pi/f^{5/3}}{4\pi/f^{5/3}} = f^{1/3} - 1
	\end{equation}
	
	\textbf{Why is this exact?}
	
	\begin{itemize}
		\item The common factor $4\pi$ cancels out
		\item The projection factor $k_{\text{geom}}$ cancels out
		\item Only the fractal exponents ($5/3$ and $4/3$) determine the ratio
		\item The result depends \textbf{only} on $f$: $f^{1/3} - 1 = 18.57$
	\end{itemize}
	
	\begin{important}{Fundamental Distinction}
		\textbf{Absolute values:}
		\begin{itemize}
			\item Depend on $k_{\text{geom}}$, $f$, and SI conversion
			\item ~2\% deviation due to higher-order quantum effects
			\item Consistent with all T0 predictions
		\end{itemize}
		
		\textbf{Ratios:}
		\begin{itemize}
			\item Depend \textbf{only} on $f$
			\item $k_{\text{geom}}$ and SI factors cancel out
			\item Mathematically exact from fractal exponents
			\item Difference $< 10^{-13}$ (numerical precision)
		\end{itemize}
		
		$\Rightarrow$ The ratio prediction is \textbf{not an approximation}, but an \textbf{exact geometric relation}!
	\end{important}
	
	\subsection{Analogy to the Koide Formula}
	
	This behavior is analogous to the Koide formula for lepton masses:
	
	\begin{itemize}
		\item \textbf{Individual masses:} ~0.4--1.2\% deviation
		\item \textbf{Koide ratio:} $\pm 0.0004\%$ precision!
	\end{itemize}
	
	The ratio is \textbf{more fundamental} than absolute values because systematic factors cancel out.
	
	\textbf{For g-2 in T0:}
	\begin{itemize}
		\item \textbf{Absolute values:} ~2\% deviation
		\item \textbf{Ratio $\Delta a(\tau-\mu)/\Delta a(\mu-e)$:} Exactly $= f^{1/3} - 1$
	\end{itemize}
	
	This is \textbf{not a weakness}, but shows the \textbf{geometric structure} of the theory!
	
	\section{Precise Ratio Predictions}
	
	\subsection{Analogy to the Koide Formula}
	
	The Koide formula for lepton masses:
	\begin{equation}
		\frac{m_e + m_\mu + m_\tau}{(\sqrt{m_e} + \sqrt{m_\mu} + \sqrt{m_\tau})^2} = \frac{2}{3} \pm 0.0004\%
	\end{equation}
	
	shows: \textbf{Ratios} are more precise than absolute values!
	
	\textbf{Question:} Does this also hold for g-2?
	
	\subsection{The Ratio of Differences}
	
	Define the differences:
	\begin{align}
		\Delta a(\mu - e) &= a_\mu - a_e = \frac{4\pi}{f^{5/3}} \\
		\Delta a(\tau - \mu) &= a_\tau - a_\mu = \frac{4\pi}{f^{4/3}} - \frac{4\pi}{f^{5/3}}
	\end{align}
	
	\textbf{Ratio:}
	\begin{align}
		\frac{\Delta a(\tau - \mu)}{\Delta a(\mu - e)} &= \frac{4\pi/f^{4/3} - 4\pi/f^{5/3}}{4\pi/f^{5/3}} \\
		&= \frac{f^{5/3}}{f^{4/3}} - 1 \\
		&= f^{5/3 - 4/3} - 1 \\
		&= f^{1/3} - 1
		\label{eq:ratio}
	\end{align}
	
	\begin{important}{Core Prediction}
		\begin{equation}
			\boxed{\frac{\Delta a(\tau - \mu)}{\Delta a(\mu - e)} = f^{1/3} - 1 = 18.57}
		\end{equation}
		
		This relation is:
		\begin{itemize}
			\item \textbf{Parameter-free} (only $f$!)
			\item \textbf{Independent} of $k_{\text{geom}}$
			\item \textbf{Exact} (difference $< 10^{-13}$)
			\item \textbf{Testable} at Belle II
		\end{itemize}
	\end{important}
	
	\subsection{Numerical Verification}
	
	With $f = 7500$:
	\begin{align}
		f^{1/3} &= 7500^{1/3} = 19.57 \\
		f^{1/3} - 1 &= 18.57
	\end{align}
	
	From T0 values:
	\begin{align}
		\Delta a(\mu - e) &= 4.373 \times 10^{-6} \\
		\Delta a(\tau - \mu) &= 8.123 \times 10^{-5} \\
		\text{Ratio} &= \frac{8.123 \times 10^{-5}}{4.373 \times 10^{-6}} = 18.57
	\end{align}
	
	\textbf{Agreement:} Perfect! ✓✓✓
	
	\subsection{Testable Prediction for Tau}
	
	With experimental values for $e$ and $\mu$:
	\begin{align}
		a_e^{\text{exp}} &= 1.160 \times 10^{-3} \\
		a_\mu^{\text{exp}} &= 1.166 \times 10^{-3} \\
		\Delta a(\mu - e)^{\text{exp}} &= 6.000 \times 10^{-6}
	\end{align}
	
	\textbf{Prediction:}
	\begin{align}
		\Delta a(\tau - \mu) &= \Delta a(\mu - e)^{\text{exp}} \times (f^{1/3} - 1) \\
		&= 6.000 \times 10^{-6} \times 18.57 \\
		&= 1.114 \times 10^{-4} \\
		a_\tau^{\text{predicted}} &= 1.166 \times 10^{-3} + 1.114 \times 10^{-4} \\
		&= 1.277 \times 10^{-3}
	\end{align}
	
	\section{The Fractal Correction $K_\text{frak}^{3/2}$: Resolving the ~2\% Deviation}
	
	\subsection{Central Insight}
	
	The ~2\% deviation of the ideal g-2 formula is \textbf{identical} to the fractal correction $K_\text{frak}^{3/2}$, already known from the mass formulas:
	
	\begin{equation}
		\boxed{k_\text{eff} = \frac{k_\text{geom}}{K_\text{frak}^{3/2}} = \frac{(2/\sqrt{\varphi})\sqrt{2}}{(1-100\xi)^{3/2}} = \frac{2.224}{0.98007} = 2.2692}
	\end{equation}
	
	The experimentally reconstructed value $k_\text{rec} = 2.2696$ confirms this to \textbf{0.014\%}.
	
	\subsection{Why the Exponent $3/2$?}
	
	The exponent $3/2 = D/2$ has a clear geometric meaning:
	\begin{itemize}
		\item The fractal correction $K_\text{frak} = 1 - 100\xi$ describes the deviation of the real lattice from the ideal crystal structure
		\item In the 4D→3D projection ($k_\text{geom}$), this correction acts across $D = 3$ spatial dimensions
		\item Per dimension: $K_\text{frak}^{1/2}$ (square root of the 1D correction)
		\item Total: $K_\text{frak}^{D/2} = K_\text{frak}^{3/2}$
	\end{itemize}
	
	\subsection{Consistency with Other T0 Corrections}
	
	\begin{table}[h]
		\centering
		\resizebox{\textwidth}{!}{\begin{tabular}{lcc}
			\toprule
			\textbf{Physical Quantity} & \textbf{$K_\text{frak}$ Exponent} & \textbf{Interpretation} \\
			\midrule
			Particle masses & $K_\text{frak}^1$ & Scalar correction (1D) \\
			g-2 projection factor & $K_\text{frak}^{3/2}$ & 3D spatial projection ($D/2$) \\
			Gravitational constant & $K_\text{frak}^1$ & Scalar correction (1D) \\
			\bottomrule
		\end{tabular}}
		\caption{Universal fractal correction with context-dependent exponent}
	\end{table}
	
	The different exponents are \textbf{not free parameters} but follow from the dimensionality of the respective physical process.
	
	\section{Experimental Tests}
	
	\subsection{Belle II (2027--2028)}
	
	Belle II expects sensitivity of $\sim 10^{-7}$ for $a_\tau$.
	
	\textbf{Test 1: Absolute value}
	\begin{itemize}
		\item T0 prediction (ideal): $a_\tau = 1.269 \times 10^{-3}$
		\item T0 prediction (corrected): $a_\tau = 1.245 \times 10^{-3}$
		\item From ratio: $a_\tau = 1.277 \times 10^{-3}$
		\item Difference: ~1\%
	\end{itemize}
	
	\textbf{Test 2: Ratio}
	\begin{itemize}
		\item T0 prediction: $\Delta a(\tau - \mu) / \Delta a(\mu - e) = 18.57$
		\item This is the \textbf{more precise} prediction!
		\item Independent of absolute calibration
	\end{itemize}
	
	\textbf{Possible outcomes:}
	\begin{enumerate}
		\item \textbf{Confirmation:} Ratio $\approx 18.6$ \\
		$\Rightarrow$ Strong evidence for fractal structure hypothesis
		
		\item \textbf{Deviation:} Ratio $\neq 18.6$ \\
		$\Rightarrow$ Different fractal dimensions or additional physics
		
		\item \textbf{Null result:} $a_\tau < 10^{-8}$ \\
		$\Rightarrow$ T0 contributions suppressed or theory needs revision
	\end{enumerate}
	
	\subsection{Fermilab/J-PARC}
	
	Further precision improvements for $a_\mu$:
	\begin{itemize}
		\item Reduction of experimental uncertainties
		\item Clearer determination of SM discrepancy
		\item Refinement of $\Delta a(\mu - e)$ measurement
	\end{itemize}
	
	\section{Comparison with Other Approaches}
	
	\begin{table}[h]
		\centering
		\resizebox{\textwidth}{!}{\begin{tabular}{lccc}
			\toprule
			\textbf{Approach} & \textbf{Precision} & \textbf{Parameters} & \textbf{Explainable} \\
			\midrule
			QED (SM) & Perfect & Many & Yes \\
			T0 (ideal, $k_\text{geom}$) & 2\% & 0 & \textbf{Completely} \\
			T0 (with $K_\text{frak}^{3/2}$) & 0.01--0.15\% & 0 & \textbf{Completely} \\
			\bottomrule
		\end{tabular}}
		\caption{Comparison of different approaches}
	\end{table}
	
	\textbf{T0 Philosophy:} We choose \textbf{explainability} over precision!
	
	\section{Proof: The Experimental Reconstruction IS the Fractal Correction}
	
	\subsection{The Central Insight}
	
	The ratio $\Delta a(\tau-\mu) / \Delta a(\mu-e) = f^{1/3} - 1$ is \textbf{mathematically exact} because the correction factor $k_{\text{geom}}$ cancels out completely.
	
	In Rev.~10, the experimentally reconstructed value $k_\text{rec} = 2.269$ was considered ``semi-empirical.'' Now it turns out: this value is \textbf{exactly the fractal correction}:
	
	\begin{equation}
		\boxed{k_\text{eff} = \frac{k_\text{geom}}{K_\text{frak}^{3/2}} = \frac{2.224}{(1-100\xi)^{3/2}} = \frac{2.224}{0.98007} = 2.2692}
	\end{equation}
	
	\subsection{Comparison: Reconstructed vs. Fractal Corrected}
	
	\begin{equation}
		k_{\text{geom}}^{\text{(reconstructed)}} = \frac{S_3/f}{a_e^{\text{(exp)}}} = \frac{2\pi^2 / 7500}{1.160 \times 10^{-3}} = 2.2696
	\end{equation}
	
	\textbf{Comparison:}
	\begin{itemize}
		\item Geometrically ideal: $k_{\text{geom}} = (2/\sqrt{\varphi}) \times \sqrt{2} = 2.224$
		\item Fractal corrected: $k_{\text{eff}} = k_{\text{geom}} / K_\text{frak}^{3/2} = 2.2692$
		\item Reconstructed from experiment: $k_{\text{geom}}^{\text{(rec)}} = 2.2696$
		\item \textbf{Difference $k_\text{eff}$ vs. $k_\text{rec}$: 0.014\%}
	\end{itemize}
	
	This proves: The ``semi-empirical'' reconstruction was all along the \textbf{purely geometric} fractal correction $K_\text{frak}^{3/2}$.
	
	\subsection{Results with Fractal Correction}
	
	\begin{table}[h]
		\centering
		\resizebox{\textwidth}{!}{\begin{tabular}{lcccc}
			\toprule
			\textbf{Lepton} & \textbf{Ideal ($k=2.224$)} & \textbf{Corr. ($k_\text{eff}=2.269$)} & \textbf{Experiment} & \textbf{Dev.} \\
			\midrule
			Electron & $1.184 \times 10^{-3}$ & $1.1598 \times 10^{-3}$ & $1.1597 \times 10^{-3}$ & \textbf{0.01\%} \\
			Muon & $1.188 \times 10^{-3}$ & $1.1642 \times 10^{-3}$ & $1.1659 \times 10^{-3}$ & \textbf{0.15\%} \\
			Tau & $1.269 \times 10^{-3}$ & $1.2454 \times 10^{-3}$ & (not measured) & Prediction \\
			\bottomrule
		\end{tabular}}
		\caption{Absolute values with geometric vs. fractal corrected $k_\text{eff}$}
	\end{table}
	
	\begin{important}{Crucial Point}
		The ~2\% deviation of the ideal formula is \textbf{fully} explained by $K_\text{frak}^{3/2}$:
		\begin{itemize}
			\item Electron: $2.03\% \to 0.014\%$ (factor 145 improvement)
			\item Muon: $1.89\% \to 0.15\%$ (factor 13 improvement)
			\item The fractal correction is \textbf{not a new parameter} but follows from $\xi = 4/30000$
		\end{itemize}
		
		Thus the T0 g-2 calculation is \textbf{fully parameter-free} with a precision of 0.01--0.15\%.
	\end{important}
	
	\subsection{Alternative: Directly from Ratio Relation}
	
	Even more precise is the calculation directly from the exact ratio:
	
	\begin{align}
		\Delta a(\mu-e)^{\text{(exp)}} &= a_\mu^{\text{(exp)}} - a_e^{\text{(exp)}} = 6.000 \times 10^{-6} \\
		\Delta a(\tau-\mu) &= \Delta a(\mu-e)^{\text{(exp)}} \times (f^{1/3} - 1) \\
		&= 6.000 \times 10^{-6} \times 18.57 = 1.114 \times 10^{-4} \\
		a_\tau^{\text{(Ratio)}} &= a_\mu^{\text{(exp)}} + \Delta a(\tau-\mu) \\
		&= 1.166 \times 10^{-3} + 1.114 \times 10^{-4} \\
		&= \boxed{1.277 \times 10^{-3}}
	\end{align}
	
	\textbf{Note:} This prediction is \textbf{independent} of $k_{\text{geom}}$ and uses only the exact geometric ratio structure!
	
	\subsection{Two Complementary Tau Predictions}
	
	\begin{table}[h]
		\centering
		\resizebox{\textwidth}{!}{\begin{tabular}{lcc}
			\toprule
			\textbf{Method} & \textbf{$a_\tau$ Prediction} & \textbf{Dependent on} \\
			\midrule
			Purely geometric & $1.269 \times 10^{-3}$ & $k_{\text{geom}} = 2.224$ (geometric) \\
			With rec. $k_{\text{geom}}$ & $1.246 \times 10^{-3}$ & $k_{\text{geom}} = 2.269$ (from $a_e$) \\
			From ratio & $1.277 \times 10^{-3}$ & Only $f$ (exact) \\
			\midrule
			Range & $1.25$--$1.28 \times 10^{-3}$ & $\pm 1.5\%$ \\
			\bottomrule
		\end{tabular}}
		\caption{Three T0 predictions for $a_\tau$}
	\end{table}
	
	\subsection{What does this mean for Belle~II?}
	
	\textbf{If Belle~II measures:}
	
	\begin{enumerate}
		\item \textbf{$a_\tau \approx 1.28 \times 10^{-3}$:}
		\begin{itemize}
			\item ✓ Confirms the exact ratio relation $f^{1/3} - 1$
			\item ✓ Shows that experimental $a_\mu$ and ratio structure are correct
			\item → \textbf{Strongest confirmation of T0 geometry}
		\end{itemize}
		
		\item \textbf{$a_\tau \approx 1.25 \times 10^{-3}$:}
		\begin{itemize}
			\item ✓ Confirms fractal correction $k_{\text{eff}} = k_\text{geom}/K_\text{frak}^{3/2}$
			\item ✓ Shows that the ~2\% deviation is fully explained by $K_\text{frak}^{3/2}$
			\item → \textbf{Strongest confirmation of fractal T0 geometry}
		\end{itemize}
		
		\item \textbf{$a_\tau \approx 1.27$--$1.28 \times 10^{-3}$:}
		\begin{itemize}
			\item ✓ Confirms exact ratio relation $f^{1/3} - 1$
			\item ? Fractal correction at tau with different exponent?
		\end{itemize}
		
		\item \textbf{$a_\tau$ outside $1.24$--$1.28$:}
		\begin{itemize}
			\item ✗ T0 structure needs revision
		\end{itemize}
	\end{enumerate}
	
	\begin{keypoint}[Key Statement]
		The ~2\% deviation of the purely geometric T0 predictions is \textbf{fully} explained by the fractal correction $K_\text{frak}^{3/2}$:
		\begin{itemize}
			\item Electron: $2.03\% \to 0.014\%$
			\item Muon: $1.89\% \to 0.15\%$
		\end{itemize}
		
		The effective projection factor $k_\text{eff} = k_\text{geom}/K_\text{frak}^{3/2} = 2.2692$ is \textbf{not a new free parameter} but follows directly from $\xi = 4/30000$.
		
		The most precise T0 prediction for tau uses the exact ratio relation:
		\begin{equation}
			\boxed{a_\tau = 1.277 \times 10^{-3}}
		\end{equation}
	\end{keypoint}
	
	\section{Important Note: No $\alpha$ in the T0 g-2 Formulas}
	
	\textbf{IMPORTANT:}
	The T0 formulas for g-2 contain \textbf{no $\alpha$}!
	
	In natural units ($\hbar = c = \alpha = 1$):
	\[ a_\ell = f(\varphi, \xi, f, \text{generation quantum numbers}) \]
	
	The anomalous moment is a \textbf{purely geometric quantity},
	following from the winding structure in the torsion lattice.
	
	Ratios like $\Delta a(\tau-\mu)/\Delta a(\mu-e) = f^{1/3} - 1$ are
	\textbf{independent} of:
	• $\alpha$ (fine-structure constant)
	• SI conversion factors
	• $k_{\text{geom}}$ (projection factor)
	
	They depend ONLY on the fractal structure!
	
	\section*{Further Reading and Resources}
	
	\textbf{T0 Theory and Python Scripts:}
	\begin{itemize}
		\item Repository: github.com/jpascher/T0-Time-Mass-Duality
		\item Python scripts: github.com/jpascher/T0-Time-Mass-Duality/blob/main/2/python/
		\item Time-Mass Duality documentation
		\item Fundamental Fractal Geometric Field Theory (FFGFT)
	\end{itemize}
	
	\textbf{Experimental Results:}
	\begin{itemize}
		\item Fermilab Muon g-2 (2025): \href{https://muon-g-2.fnal.gov/}{muon-g-2.fnal.gov}
		\item Theory Initiative White Paper
		\item Belle II: \href{https://www.belle2.org/}{www.belle2.org}
	\end{itemize}
	
	\textbf{Related T0 Documents:}
	\begin{itemize}
		\item Lepton masses: Systematic derivation from quantum numbers
		\item Koide formula in T0: Geometric interpretation
		\item Fractal spacetime: $D_f = 3 - \xi$
	\end{itemize}
	